\section{怎样编写作文提纲}

每写一篇作文，总要先审题，再构思，然后在审题、构思的基础上写出文章的纲要，这就叫编写作文提纲。开始下笔作文之前，如果认真地编好提纲，就可预见到文章写出后大体是什么样子，做到胸有成竹；不然的话，很容易造成条理不清，详略失调，或遗漏了某些重要内容，或出现了多余的部分，这一段用什么方法写或那一段用什么方法写都将心中无数。与其错了再改（如是考试，即使发现错了，也没有多少时间改），何不事先绘好蓝图，搭好架子呢？

要想编好作文提纲，应该注意以下几个方面。

\textbf{1. 要善于落实构思成果，把文章中心突现出来。}

编写提纲的第一步工作，就是根据审题和构思而明确起来的中心意思，用一两句话简明扼要地写出来，使要写的作文有一个统帅全文的中心思想，从而表明自己对某种人、某件事、某种现象或某种道理的根本态度和认识。这是很重要的，因为一篇文章如果没有中心，就好象一个人没有灵魂一样。

例如有位同学为作文题《我们的张老师》编写了这样一份提纲：

①\ 概括张老师给我们的初步印象。

②\ 描写张老师的外貌和几个口头禅。

③\ 具体描写张老师的两件事。

④\ 用张老师批评我们给他起外号结束全文。

可以看出，这位同学虽已编了提纲，但是文章的中心意思是什么，恐怕连他自己都不清楚，写出的文章怎么可能中心明确又有积极意义呢！有的同学甚至把编写作文提纲当作例行公事，应付一下，并不把提纲当成作文的依据，使得提纲是提纲，文章是文章，彼此对不上号，那就更不好了。

我们看另一位同学为《我们的张老师》编写的提纲：

\noindent
[作文提纲]

一、中心思想：通过张老师既严格要求又无比关怀我们同学的动人事例，歌颂人民教师美好的心灵。

二、段意和写法：

（1）描写初次见面的第一堂课，他两次叫我们起立敬礼的情景。（略写印象 --- “严”）

（2）记叙张老师令人难忘的三件事。（详写对我们的“严”与“爱”）

①\ 有一次，他叫二十多个同学将他批改的作文一笔不差地誊抄出来，再给他批改；改过之后，又叫三个把字写得东倒西歪的同学再工整地抄给他看。至今全班同学作业本没有潦草字。

②\ 曹诚同学在运动会短跑中不幸跌倒了，摔伤了腿，他跑去借了小平车，亲自拉曹诚同学上医院，我们要拉他不肯，感动得曹诚都哭了。

③\ 有一次他病了，我们几个班干部代表同学们去看望他，哪知他趴在床上还在给我们编写复习讲义呢。

（3）抒发对我们张老师热爱的感情，他真是我们的好老师。（用抒情笔调结尾，文字要简洁）。

这份提纲由于中心明确，善于从构思的素材中提炼出一个有积极意义的中心意思，使得作文有了“纲”，选材也就有所遵循了。

\textbf{2. 要善于根据中心思想，列出切实的内容条目，并使条目之间有严密的逻辑联系。}

如果把中心思想比作“纲”的话，那么所要写的具体内容就是“目”。符合要求的作文提纲应该做到纲举目张。光有目而无纲，文章就失去统帅全篇的灵魂；光有纲而无目，文章就失去骨骼和血肉，这都是不符合要求的。只有在确定了中心意思之后，再把表现这一中心意思的材料确定下来，行文才可能思路畅达，条理清楚，得心应手。例如上面所列举的《我们的张老师》第二份作文提纲，就有纲举目张的特色，它善于围绕中心思想进行选材，把由表面严厉而严肃中透爱、爱是本质的思路条理清楚地反映了出来，而且把表现的手法也规定了下来，这样再动笔作文，当然就水到渠成了。

我们再看看以《说忧乐》为题而编写的两份提纲：

[例一]

中心思想：论述只有以人民的忧乐为忧乐，才是革命者的忧乐观。

段落大意：

（1）从自古以来就有两种忧乐观谈起。

（2）论述革命者的忧乐观。

（3）批驳个人主义者忧乐观。

（4）结论：我们要以人民的忧乐为忧乐。

[例二]

中心思想：同[例一]

段意与写法：

（1）从范仲淹的名言谈起，用大汉奸秦桧为了自己的私利不惜杀害忠良、出卖祖国进行对比，说明自古以来，就有两种完全对立的忧乐观。（略写）

（2）着重论证：在无产阶级革命时代的共产主义战士一心为党和人民的命运担忧，以谋取人民幸福为己任，这才是当代最崇高的忧乐观。（方志敏为了“可爱的中国”慷慨献身；叶挺为了人民的解放“愿把牢底坐穿”；董存瑞为了解放全中国而舍身歼敌；彭老总要为人民鼓与呼，屡遭迫害，不屈至死；张志新直面鬼魅发出真理的呐喊，英勇牺牲······）

（3）深入分析：这些革命战士之所以如此，因为他们心中只有人民的忧乐，无私才能无畏。他们是民族的脊梁，我们的楷模。

（4）进行对比：斥责那些终日为权力的大小忧、为待遇的高低忧、为儿孙的工作忧、为子女的婚事忧，为饱食终日乐、为中饱私囊乐、为大开后门乐、为损公肥私乐的人。（略写）

（5）时代需要我们以祖国、人民的忧乐为忧乐，我们要做新时代大写的人。（略写作结）

比较这两份作文提纲，前一份失之过简，尽管中心明确，层次分明，却没能反映出该写些什么和用什么手法写。作者动起笔来，还得边写边想材料，很可能造成堆砌材料或者没有材料，材料不能证明论点，遗漏重要材料等情况，在

写法上也心中无数，哪儿先写，哪儿后写，哪儿该详，哪儿该略，用什么论证手法证明，段与段之间该如何环环相扣又层层深入，等等，都没有把握安排妥当。这大约是在许多初学写作者通常惯用的提纲类型。后一份则不然，它不仅按照表达中心的需要确定了所选的具体材料，而且考虑了详略、主次，设计了各种论证方法，注意到各个条目之间的逻辑联系，甚至文中该用的一些重点词句都列了出来，可以想见，按这份提纲落笔成文，当然就能做到下笔千言、倚马可待了。

\textbf{3. 必须以纲带目，纲目结合，决不可以目漫编，纲目分离。}

有一些作文题，本身并不能直接显示出中心意思，例如记叙文题目《生动的一课》、《闲不住的人》、《暑假记事》、《在大路上》、《激动人心的一刻》、《他笑了》、《我的第一位班主任》、《毕业前夕》、《春游花果山》……，议论文题目《从拔河想到的》、《一滴水的启示》、《说“难”》、《羡慕小议》、《谈青年时代》、《我们应怎样渡过青春》、《评比上不足、比下有余》······这一类作文题，可以根据自己的经历、感受和认识，确定不同的中心意思。“纲”不同，“目”随之也会不同，所以必须要以纲带目，纲目吻合，而不能相反。如果先将杂乱的材料和想法罗列出来，然后再勉强凑个中心意思，那就把“纲”与“目”的关系弄颠倒了，结果“目”中所列的材料与“纲”不相吻合，甚至会搞得驴唇不对马嘴。例如有这样一份提纲：

作文题：一滴水的启示

作文提纲：

（1）中心论点：要想成材，必须把自己汇入四化洪流中去。

（2）论据：

①\ 一滴水不可忽视，它能反映整个世界，我们不能小看自己的作用。

②\ 一滴水是渺小的，一阵风就能把它吹没。我们每个人的力量也是渺小的，只要脱离了集体，就会茫然不知所向。

③\ 一滴水怎样才能不干涸？释迦牟尼说：“把它放到大海去。”我们要想成材，只有把自己汇入四化洪流中去。

④\ 浩瀚的大海包含着无数的一滴水，所以“一滴水”是可贵的。我们每个青年都是实现四化的有用人才，要特别珍视自己的才华。

⑤\ 许许多多的一滴水汇合到一起，就能形成了不起的力量。我们青年一代团结起来共同奋斗，就能成为四化建设的浩荡大军。

这是一位正在进修的社会青年所写的提纲。它的优点是中心明确又有积极的现实意义，并能恰当、熟练地运用类比论证方法。但是，它也存在着明显的毛病，那就是没有以“纲”带“目”，而是以“目”凑“纲”，摆出了一些杂乱的论据，却不管这些论据能否去证明中心论点，也没有考虑到前后层次之间的严密的逻辑联系。如果把“要想成材，必须把自己汇入四化洪流中去”作为中心论点（这个论点本身是正确的），那么，对上述论据应该进行修改：

首先，将①删去，因为它与中心论点无关。

其次，将②与⑤合并，提法应改变为：一滴水是渺小的，许许多多的一滴水汇合到一起才能形成伟大的力量；我们每个人同样是渺小的，只有一起为四化建设献身，才能形成一支浩荡的大军。

再次，要改变④的角度，不是强调一滴水的可贵，而是强调无论多么闪光的一滴水，一旦离开了伟大的集体，就会失去它应有的作用。

最后，要将③强调出来，因为它正是自己着力论证的中心。可以把③放在全文的首段，提出问题，点明中心；然后再围绕这个中心展开论证。

\textbf{4. 坚持做拟定作文提纲或修订作文提纲的练习，仔细推敲，认真比较，就能打开思路，熟能生巧，练出拟好提纲的能力，为写好作文打下坚实的基础。}

作文提纲的基本要求是：中心突出，纲举目张，要点明确，层次分明，有主有次，详略得当，组织有序，结构谨严，语言既不要失之太简，也不要巨细不捐，显得冗长。我们根据这些要求，作一些拟定或修改作文提纲的练习，必将有明显的效果。

例如，有位同学曾有意识地作过拟写作文提纲的练习，开始时失之过简，后来想得面宽了，而往往使所用的材料与中心思想不相吻合，有时条目之间缺乏内在的逻辑联系，但他锲而不舍，越练越好，终于能写出象样的提纲，作文有了显著的进步。我们试引他的两份作文提纲，比较一下看看，究竟哪一份提纲较好，好在哪里，怎样学习，哪一份提纲较差，差在何处，怎样修订？请你思考。

[\ 1\ ]\ 作文题：失败是成功之母

（1）中心论点：善于总结失败的教训，并能坚持下去，终能成功。

（2）段意与写法

①\ 提出面对失败或挫折应该怎么办，引出中心论点。（略写）

②\ 运用对比法进行反证（\uline{张国焘叛变开小差}；\\ 
\uline{庞涓骄傲轻敌自刎马陵道}；有的同学因为期中考试成绩不好就灰心丧气，成绩日下。）

③\ 运用引证法进行正面论证（荀子：“锲而不舍，金石可镂。”
\uline{孟子：“生于忧患，死于安乐”}。陈毅同志说：“一喜有错误，痛改便光明。”）

④\ 运用例证法进行正面论证（欧立希发明药物六○六，失败了六百零五次，坚持实验，终于成功；\\ 
\uline{达·芬奇画蛋不止，终成一代宗师}；\\ 
陈景润演算几麻袋草稿纸，终于攻下了“哥德巴赫猜想”；
\uline{曹雪芹十载辛苦，终于写出了震古烁今的《红楼梦》}。）

⑤\ 分析失败与成功的辩证关系，点题。（略写）

⑥\ 结论，呼应开头，深化中心思想。（略写）
\newpage

[\ 2\ ]\ 作文题：求知无坦途

（1）中心论点：要想求得真知，必须不畏险阻，迎难而上，企图走坦途则将一事无成。

（2）结构提纲：

第一，引论：从求知过程中碰到种种困难，提出问题，点明中心论点。（略写）

第二，本论：

①\ 事实论证（“唐宋八大家”之一的苏洵不怕晚学，尽焚其稿，终于成名；华罗庚不顾家穷、体病、学历浅而顽强奋进，成为闻名全球的数学家；张海迪无视严重残疾，刻苦攻读，译出名著，成为“青年楷模”。详写。）

②\ 理论论证（马克思的名言，王安石《游褒禅山记》中的警句。）

③\ 对比论证（只图轻松者和知难而上者，半途而废者和坚持不懈者；失败者和成功者。略写。）

④\ 进层论证（比较容易求得的是浅知，难于求得的才是真知。）

第三，结论：一切有志求取真知者，必须无畏险阻，沿着崎岖的山路奋勇攀登，终究可以达到光辉的顶峰。

这两份提纲内容都较充实，中心都颇明确。但是第一份提纲中的有些材料并不能证明中心论点，前后层次的连贯性也不够好。比较起来，第二份提纲就成熟多了。你的看法如何呢？
\newpage